\documentclass[12pt, a4paper]{letter}
\usepackage[utf8]{inputenc}
\usepackage[T1]{fontenc}
\usepackage[english]{babel}
\usepackage[top=2.5cm, bottom=2.5cm, left=3cm, right=3cm]{geometry}
\usepackage[version=4]{mhchem}

\signature{Federico Garc\'ia Cresp\'i\\
           Departamento de Ingenier\'ia de Computadores\\
           Universidad Miguel Hern\'andez de Elche\\
           Elche (Alicante), Spain\\
           \texttt{fedeg@umh.es}\\
           ORCID: 0000-0002-7129-7791}

\address{Federico Garc\'ia Cresp\'i\\
         Departamento de Ingenier\'ia de Computadores\\
         Universidad Miguel Hern\'andez de Elche\\
         Av. de la Universidad s/n\\
         03202 Elche, Alicante, Spain}

\begin{document}

\begin{letter}{The Editor-in-Chief\\
               Air Quality, Atmosphere \& Health\\
               Springer Nature}

\opening{Dear Editor,}

We submit for your consideration the manuscript entitled
\textit{``Operational predictability limits for urban \ce{PM10} and
\ce{PM_{2.5}} air-quality forecasting in Valencia and Madrid''},
which we believe is well suited to
\textit{Air Quality, Atmosphere \& Health}.

\textbf{The problem.}
Short-term forecasting of particulate matter is relevant not only for
model comparison, but also for episode management, public-health
alerts, and operational decision support in urban air-quality systems.
However, much of the published forecasting literature still relies on
aggregate error metrics from static hold-out evaluations, without
comparing performance against a persistence baseline horizon by horizon
and without ensuring train-only preprocessing. Under those conditions,
it is difficult to determine whether an apparently accurate model
provides genuine operational added value at the lead times that matter
for warning and response.

\textbf{What we do.}
This manuscript quantifies the operational predictability limit $H^*$
for daily \ce{PM10} and \ce{PM_{2.5}} concentrations across 43 valid
station-level series from Valencia and Madrid. All models are evaluated
under a strict rolling-origin, expanding-window protocol with
train-only preprocessing. We compare simple persistence, seasonal
persistence, SARIMA, and XGBoost, and summarize horizon-by-horizon
skill relative to persistence using two complementary descriptors:
$H^*_{\mathrm{relax}}$ (farthest horizon with positive skill) and
$H^*_{\mathrm{strict}}$ (longest contiguous block of positive skill).

\textbf{Key findings.}
Simple persistence produced $H^*_{\mathrm{relax}} =
H^*_{\mathrm{strict}} = 0$ in all 43 evaluated series, validating the
zero-skill reference. SARIMA achieved $H^*_{\mathrm{relax}} = 7$~days
in all series, while XGBoost reached that horizon in the large
majority of series. Across models, Valencia was generally more
predictable than Madrid, and \ce{PM_{2.5}} showed more continuous skill
than \ce{PM10}. The central practical result is that
$H^*_{\mathrm{strict}}$, not only $H^*_{\mathrm{relax}}$, is the
relevant descriptor for communicating consecutive-day warning lead
times and judging whether a forecasting system is useful for deployment.

\textbf{Fit with Air Quality, Atmosphere \& Health.}
We believe the manuscript fits the journal because it links air-quality
forecasting directly to public-health relevance and operational use.
Rather than presenting methodology in isolation, the paper provides a
city- and station-specific framework for deciding at which forecast
horizons advanced models offer meaningful value over persistence in
real urban monitoring settings. This is directly relevant to alert
systems, episode management, and deployment decisions in air-quality
practice.

\textbf{Declarations.}
This manuscript has not been submitted elsewhere. A preprint of a
related methodological paper is available on arXiv
(arXiv:2603.20315); the present manuscript provides a new
multi-station, multi-city, multi-pollutant empirical analysis with new
results. All code is publicly available at\\
\texttt{https://github.com/fedeg-umh-es/Hstar\_PM10\_PM25\_Madrid\_Valencia}.
The author declares no competing interests.

\closing{Sincerely,}

\end{letter}
\end{document}
